\documentclass[11pt]{article}
\usepackage[a4paper,margin=1in]{geometry}
\usepackage{amsmath,amssymb,amsthm}
\usepackage[T1]{fontenc}
\usepackage{lmodern}
\usepackage[hidelinks]{hyperref}
\setlength{\emergencystretch}{3em}

\theoremstyle{plain}
\newtheorem{theorem}{Theorem}
\newtheorem{lemma}[theorem]{Lemma}
\theoremstyle{remark}
\newtheorem{remark}[theorem]{Remark}

\title{A Proof of Conjecture 00000000034}
\author{%
  Feiyu\thanks{Independent researcher.}
  \and
  Claude (Opus 5)\thanks{Anthropic.}%
}
\date{2026-09-12}

\begin{document}
\maketitle

\begin{abstract}
Every finite colouring of $\mathbb N$ contains a monochromatic three-term
geometric progression $x,\,xr,\,xr^2$ with integer ratio $r>1$. The proof is
the standard reduction to van der Waerden's theorem along the powers of two.
The argument is formalized in Lean~4 against Mathlib, and the formalization is
unconditional: van der Waerden's theorem is derived from Mathlib's
Hales--Jewett theorem rather than assumed.
\end{abstract}

\section{Statement}

\begin{theorem}\label{thm:main}
Let $c:\mathbb N\to K$ be any colouring of the natural numbers by a finite set
$K$ of colours. Then there exist a positive integer $x$ and an integer $r>1$
with
\[
  c(x)=c(xr)=c(xr^2).
\]
\end{theorem}

Since the witness produced below satisfies $x\ge1$, it is immaterial whether
$\mathbb N$ is taken to contain $0$.

\section{Proof}

We use van der Waerden's theorem in the following form.

\begin{lemma}[van der Waerden, three-term case]\label{lem:vdw}
For every colouring $c':\mathbb N\to K$ by a finite set of colours there exist
$a\ge0$ and $d>0$ with
\[
  c'(a)=c'(a+d)=c'(a+2d).
\]
\end{lemma}

\begin{proof}[Proof of Theorem~\ref{thm:main}]
Define a colouring of the exponents by
\[
  c'(m)=c(2^m),\qquad m\in\mathbb N .
\]
This is again a colouring by the finite set $K$, so Lemma~\ref{lem:vdw} gives
$a\ge0$ and $d>0$ with
\[
  c(2^a)=c(2^{a+d})=c(2^{a+2d}).
\]
Put
\[
  x=2^a,\qquad r=2^d .
\]
Then $x\ge1$, and $r=2^d>1$ because $d>0$. The exponent laws give
\[
  2^{a+d}=2^a\,2^d=xr,
  \qquad
  2^{a+2d}=2^a\left(2^d\right)^2=xr^2 ,
\]
so the displayed equalities read $c(x)=c(xr)=c(xr^2)$. As $r$ is an integer
greater than~$1$, this is a monochromatic geometric progression of the required
form.
\end{proof}

\begin{remark}
Nothing is special about the base $2$: any integer base $b\ge2$ works, and
applying van der Waerden's theorem to $k$-term arithmetic progressions instead
yields monochromatic geometric progressions $x,xr,\ldots,xr^{k-1}$ of every
length $k$. Only the three-term case is formalized below.
\end{remark}

\section{Formal verification}

The accompanying Lean~4 project formalizes Theorem~\ref{thm:main} against
Mathlib. The formalization is \emph{unconditional}: van der Waerden's theorem
is not taken as a hypothesis.

\begin{itemize}
  \item \texttt{vanDerWaerden3} --- Lemma~\ref{lem:vdw}, obtained from Mathlib's
    \texttt{Combinatorics.exists\_mono\_homothetic\_copy} applied to the finite
    set $S=\{0,1,2\}\subseteq\mathbb N$. That theorem --- a corollary of the
    Hales--Jewett theorem, which Mathlib proves --- produces $a>0$, $b$ and a
    colour $k$ with $c'(a\cdot s+b)=k$ for every $s\in S$, which is exactly a
    monochromatic three-term arithmetic progression with first term $b$ and
    common difference $a$.
  \item \texttt{exists\_mono\_geometric\_progression} ---
    Theorem~\ref{thm:main} for an arbitrary finite colour type:
    \[
      \forall\,\kappa\ \text{finite},\ \forall\,c:\mathbb N\to\kappa,\quad
      \exists\,x\,r,\ 0<x\ \wedge\ 1<r\ \wedge\
      c\,x=c\,(xr)\ \wedge\ c\,(xr)=c\,(xr^2).
    \]
  \item \texttt{exists\_mono\_geometric\_progression\_fin} --- the same
    statement for a colouring $\mathbb N\to\mathrm{Fin}\,q$, matching the
    ``$q$ colours'' phrasing of the conjecture.
\end{itemize}

The project builds with \texttt{lake build} on
\texttt{leanprover/lean4:v4.33.1} with Mathlib \texttt{v4.33.1}. It contains no
\texttt{sorry}, no \texttt{admit} and no \texttt{native\_decide}, and
introduces no axioms of its own:
\texttt{\#print axioms exists\_mono\_geometric\_progression} reports exactly the
three standard axioms \texttt{propext}, \texttt{Classical.choice} and
\texttt{Quot.sound}.

\end{document}
